\documentclass[10pt,a4paper]{amsart}
\usepackage[margin=19mm]{geometry}
\usepackage{amsmath,amssymb,mathtools}
\usepackage[scr=rsfs,cal=euler]{mathalfa}
\usepackage{xfrac}
\usepackage[unicode,hidelinks,bookmarks=false]{hyperref}
\newcommand{\ie}{i.e.\ }
\newcommand{\cf}{cf.\ }
\newcommand{\resp}{respectively\ }
\newcommand{\Subsection}{\subsection}
\providecommand{\nobreakdash}{\nobreak\hbox{-}}
\setlength{\emergencystretch}{2em}
\multlinegap0pt
\usepackage{bm}
\usepackage{bbm}
\usepackage{dsfont}
\usepackage{xspace}
\usepackage{cancel}
\usepackage{mathtools}
\usepackage{amsthm}
\makeatletter
\@ifundefined{theorem}{\newtheorem{theorem}{Theorem}}{}
\makeatother
\DeclareMathOperator{\Aut}{Aut}
\DeclareMathOperator{\VCSP}{VCSP}
\newcommand{\bB}{{\mathfrak B}}
\newcommand{\bC}{{\mathfrak C}}
\title[The Complexity of Resilience Problems via Valued Constraint ]{The Complexity of Resilience Problems via Valued Constraint Satisfaction}
\author{Manuel Bodirsky, \v{Z}aneta Semani\v{s}inov\'a, Carsten Lutz}
\date{Source: arXiv:2309.15654v6 (CC BY 4.0)}
\begin{document}
\maketitle

\noindent\emph{Source and licence.} The Complexity of Resilience Problems via Valued Constraint Satisfaction. Version arXiv:2309.15654v6 (CC BY 4.0), distributed under the Creative Commons Attribution 4.0 International licence, as recorded in the arXiv licence metadata for this exact version and shown on the abstract page https://arxiv.org/abs/2309.15654v6. Full rights notice: licenses/arxiv-2309.15654-CC-BY-4.0.txt.\par\smallskip

\noindent\textit{Editorial scope.} Counted unit: the Theorem whose source label is \texttt{thm:red-fin} in the article (source0.tex lines 1455--1464, source label \texttt{thm:red-fin}), delivered together with the article's own complete original proof of it (source0.tex lines 1466--1585, one \texttt{proof} environment of 120 lines). No mathematics is written, repaired or re-derived: statement and proof are the article's own text line for line. Required context retained uncounted: none. Every definition, display and label of those lines is retained unchanged. Omitted from this package and not redistributed: 1--1454, 1465--1465, 1586--3137. Nothing inside a retained excerpt is omitted or altered: every retained body is delivered exactly as the article's lines are written, so no omission receipt of any kind accompanies this package. Citations: the retained excerpts carry no citation command at all, so no bibliography entry of the article is transcribed and no reference is redistributed. Reference adaptation: none was needed and none was made; every reference inside the delivered unit resolves inside it, so no unresolved reference or citation is produced. Printed slips kept literal: no printed word, symbol, hypothesis or number of the counted statement or of its proof was altered. Layout and class: the article's own class is \texttt{acmart} with packages including inputenc, amsmath, amsxtra, amsfonts, amssymb, hyperref, xspace, geometry, framed, todonotes, enumerate; the wrapper is the lane's pinned \texttt{amsart} editorial wrapper at 10pt on A4, which restates the theorem-like environments the retained text needs and adds the article's own macro definitions that the retained text needs (\texttt{\textbackslash Aut}, \texttt{\textbackslash VCSP}, \texttt{\textbackslash bB}, \texttt{\textbackslash bC}), with extra wrapper packages (bm, bbm, dsfont, xspace, cancel, mathtools). The delivered numbering is the unit's own. Assets: no retained span contains a figure environment, a graphics-inclusion command, a PDF-page inclusion, a table or a TikZ picture; no image file is copied into this package, the package declares no asset dependency and no image file is redistributed with it. Licence: version arXiv:2309.15654v6 of this article is distributed under the Creative Commons Attribution 4.0 International licence, as recorded in the arXiv licence metadata for this exact version and shown on the abstract page https://arxiv.org/abs/2309.15654v6. A full rights notice is delivered at licenses/arxiv-2309.15654-CC-BY-4.0.txt. Characters: every retained excerpt is pure ASCII, so the delivered engine, XeLaTeX with its default Latin Modern text font, reports no missing character.
\begin{theorem}\label{thm:red-fin}
Let $\Gamma$ be a valued structure with a finite signature such that $\Aut(\Gamma)$ contains the automorphism group of a finitely bounded
homogeneous structure $\bB$.
Let $r$ be the maximal arity of the relations of $\bB$ and the valued relations in $\Gamma$,
let $v$ be the maximal number of variables that appear in a single conjunct
of the universal sentence $\psi$ that describes the age 
of $\bB$,
and let $m \geq \max(r+1,v,3)$. 
Then there is a polynomial-time reduction from $\VCSP(\Gamma)$ to $\VCSP(\Gamma^*_{\bB,m})$.
\end{theorem}

\begin{proof}
Let $\tau,\tau^*,\sigma$ be the signatures of $\Gamma$, $\Gamma^*_{\bB,m}$, and $\bB$, respectively.  Let $\phi$ be an instance of 
$\VCSP(\Gamma)$ with threshold $u$ and let $V$ be the variables of $\phi$.  
Create a variable $y({\bar x})$ for every 
$\bar x =(x_1, \dots, x_m) \in V^m$. For every summand $R(x_1,\dots,x_k)$ of $\phi$ and
we create a summand $R^*(y(x_1, \dots, x_k, \dots, x_k))$; this makes sense since $m\geq r$. 
For every $\bar{x}, \bar{x}' \in V^m$, $p \in \{1,\dots,m\}$, and 
$i,j \colon \{1,\dots,p\} \to \{1,\dots,m\}$, add the summand 
$C_{i,j}(y(\bar x),y(\bar x'))$
if $(x_{i(1)},\dots, x_{i(p)}) = (x'_{j(1)},\dots, x'_{j(p)})$; we will refer to these as \emph{compatibility constraints}.  
Let $\phi^*$ be the resulting $\tau^*$-expression. 
Clearly, $\phi^*$ can be computed from $\phi$ in polynomial time. 


Suppose first that $(\phi,u)$ has a solution;
it will be notationally convenient to view the solution as a function $f$ from the variables of $\phi$ to the elements of $\Gamma$ (rather than a tuple). We claim that the map $f^*$ which maps $y(\bar x)$ to the orbit of $f(\bar x)$ in $\Aut(\bB)$ is a solution for $(\phi^*,u)$. And indeed, 
each of the summands involving a symbol $C_{i,j}$ evaluates to $0$,
and 
$(\phi^*)^{\Gamma^*_{\bB,m}}$ equals 
$\phi^{\Gamma}$. 


Now suppose that $(\phi^*,u)$ has a solution $f^*$. 
To construct a solution $f$ to $(\phi,u)$, we first define an equivalence relation $\sim$
on $V$. For $x_1,x_2 \in V$, define $x_1 \sim x_2$
if a (equivalently: every) tuple $t$ in $f^*(y(x_1,x_2,\dots,x_2))$
satisfies $t_1=t_2$.
Clearly, $\sim$ is reflexive and symmetric.
To verify that $\sim$ is transitive, suppose that 
 $x_1 \sim x_2$ and $x_2 \sim x_3$. 
 In the following we use that $m \geq 3$. 
Let $i$ be the identity map on $\{1,2\}$, 
let $j \colon \{1,2\} \to \{2,3\}$ be given by $x \mapsto x+1$, 
and let $j' \colon \{1,2\} \to \{1,3\}$ be given by $j'(1) = 1$ and $j'(2) = 3$. Then $\phi^*$ contains the conjuncts 
\begin{align*}
    & C_{i, i}(y(x_1, x_2, x_2, \dots, x_2), y(x_1,x_2,x_3,\dots,x_3)),\\ & C_{i, j}(y(x_2,x_3,x_3,\dots,x_3),y(x_1,x_2,x_3,\dots,x_3)), \\ 
    & C_{i, j'}(y(x_1,x_3,x_3,\dots,x_3),y(x_1,x_2,x_3,\dots,x_3)). 
\end{align*}
Let $t$ be a tuple from $f^*(y(x_1,x_2,x_3,\dots,x_3))$. 
Then it follows from the conjuncts with the relation symbols $C_{i,i}$ and $C_{i,j}$ that $t_1=t_2$ and $t_2=t_3$, and therefore $t_1=t_3$. Thus we obtain from the conjunct with $C_{i, j'}$ that $x_1 \sim x_3$. 

%\medskip 
\paragraph{Claim 0.} For all equivalence classes $[x_1]_\sim,\dots,[x_m]_\sim$, tuple $t \in f^*(y(x_1,\dots,x_m))$, $S \in \sigma$ of arity $k$, and a map $j \colon \{1, \dots, k\} \rightarrow \{1, \dots, m\}$,
whether 
$\bB \models S(t_{j(1)},\dots,t_{j(k)})$ does not depend on the choice of the representatives $x_{1},\dots,x_{m}$. 
It suffices to show this statement if we choose another representative $x'_i$ for $[x_i]_\sim$ for some $i \in \{1, \dots, m\}$, because the general case then follows by induction. 

Suppose that for every $t \in f^*(y(x_1,\dots,x_m))$ we have $\bB \models S(t_{j(1)},\dots,t_{j(k)})$ and we have to show that for every tuple $t' \in f^*(y(x_1,\dots,x_{i-1},x_i',x_{i+1},\dots,x_m))$ 
holds $\bB \models S(t'_{j(1)},\dots,t'_{j(k)})$. 
If $i \notin \{j(1),\dots,j(k)\}$, 
then $\phi^*$ contains $$C_{j,j}(y(x_1,\dots,x_m),y(x_1,\dots,x_{i-1},x_i',x_{i+1},\dots,x_m))$$
and hence 
$\bB \models S(t'_{j(1)},\dots,t'_{j(k)})$. 
Suppose $i \in \{j(1),\dots,j(k)\}$; for the sake of notation we suppose that $i=j(1)$. 
By the definition of $\sim$, and since $x_{j(1)} \sim x'_{j(1)}$, every 
 $t'' \in f^*(y(x_{j(1)},x_{j(1)}',\dots,x_{j(1)}'))$ satisfies $t''_1 = t''_2$. 
Let $\tilde t$ be a tuple from 
$$f^*(y(x_{j(1)},\dots,x_{j(k)},x'_{j(1)},\dots,x'_{j(1)})).$$  
(Here we use that $m \geq r+1$.)
\begin{itemize}
    \item $\bB \models S(\tilde t_{1},\dots,\tilde t_{k})$, because we have a compatibility constraint in $\phi^*$ between variables
$y(x_1,\dots,x_m)$ and $y(x_{j(1)}$,
$\dots,x_{j(k)},x'_{j(1)},\dots,x'_{j(1)})$; 
\item $\tilde t_1 = \tilde t_{k+1}$ because of $x_{j(1)}\sim x_{j(1)}'$  and a compatibility constraint between 
$y(x_{j(1)},\dots,x_{j(k)},x'_{j(1)},\dots,x'_{j(1)})$ and $y(x_{j(1)},x_{j(1)}',\dots,x_{j(1)}')$ 
in $\phi^*$; 
\item hence, $\bB \models S(\tilde t_{k+1},\tilde t_2,\dots,\tilde t_k)$; 
\item $\bB \models S(t'_{j(1)},t'_{j(2)},\dots,t'_{j(k)})$ 
due to a compatibility constraint between $y(x_{j(1)},\dots,x_{j(k)}$,
$x'_{j(1)},\dots,x'_{j(1)})$ and $y(x_1,\dots,x_{i-1},x'_i,x_{i+1},\dots,x_m)$
in $\phi^*$:
namely, consider the map $j' \colon \{1, \dots, k\} \to \{1, \dots, m \}$ that coincides with the identity map except that $j'(1) := k+1$, then $
\phi^*$ contains 
\begin{align*}
    C_{j,j'} \big (& y(x_1,\dots,x_{i-1},x'_{i},x_{i+1},\dots,x_m), y(x_{j(1)},\dots,x_{j(k)},x'_{j(1)},\dots,x'_{j(1)}) \big ).
\end{align*}
\end{itemize}

This concludes the proof of Claim 0. 

\medskip 

Now we can define a structure $\bC$ in the signature $\sigma$
on the equivalence classes of $\sim$. 
If $S \in \sigma$ has arity $k$, $j_1,\dots,j_k \in \{1,\dots,m\}$, and 
$[x_1]_\sim,\dots,[x_{m}]_\sim$ are equivalence classes of $\sim$ such that the tuples $t$ in $f^*(y(x_1,\dots,x_m))$ satisfy $S^{\bB}(t_{j_1},\dots,t_{j_k})$ for some representatives $x_1, \dots x_m$ (equivalently, for all representatives, by Claim 0), then add $([x_{j_1}]_\sim,\dots,[x_{j_k}]_\sim)$ to $S^{\bC}$.
No other tuples are contained in the relations of $\bC$.

%\medskip 
\paragraph{Claim 1.} 
If $[x_1]_\sim,\dots,[x_m]_\sim$ are equivalence classes of $\sim$, and 
$t \in f^*(y(x_1,\dots,x_m))$, then $[x_i]_{\sim} \mapsto t_i$,
for $i \in \{1,\dots,m\}$, is an isomorphism
between a substructure of $\bC$ and a substructure of $\bB$ for any choice of representatives $x_1, \dots, x_m$. 
First note that $[x_i]_{\sim} = [x_j]_{\sim}$ if and only if $t_i = t_j$, so the map is well-defined and bijective. 
Let $S \in \sigma$ be of arity $k$ and 
$j \colon \{1, \dots, k\} \to \{1, \dots, m \}$. 
If $\bB \models S(t_{j(1)},\dots,t_{j(k)})$, 
then 
$\bC \models S([x_{j(1)}]_\sim,\dots,[x_{j(k)}]_\sim)$ by the definition of $\bC$. 
Conversely, suppose that 
$\bC \models S([x_{j(1)}]_\sim,\dots[x_{j(k)}]_\sim)$. 
By Claim 0 and the definition of $\bC$, there is $t' \in f^*(y(x_1,\dots,x_m))$ such that $\bB \models S(t'_{j(1)},\dots,t'_{j(k)})$. Since $f^*(y(x_1,\dots,x_m))$ is an orbit of $\Aut(\bB)$, we have $\bB \models S(t_{j(1)},\dots,t_{j(k)})$ as well.

%\medskip 
\paragraph{Claim 2.} $\bC$ embeds into $\bB$. 
It suffices to verify that $\bC$ satisfies each conjunct of the universal sentence $\psi$. 
Let $\psi'(x_1,\dots,x_q)$ be such a conjunct, and let $[c_1]_\sim,\dots,[c_q]_\sim$ be elements of $\bC$. Consider the orbit $f^*(y(c_1,\dots,c_q,\dots,c_q))$ of $\Aut(\Gamma)$; this makes sense since $m \geq v$. Let $t \in f^*(y(c_1,\dots,c_q,\dots,c_q))$.
Since $t_1, \dots, t_q$ are elements of $\bB$, the tuple $(t_1,\dots,t_q)$ satisfies $\psi'$.
Claim 1 then implies that $([c_1]_\sim,\dots,[c_q]_\sim)$ satisfies $\psi'$. 

\medskip

Let $e$ be an embedding of $\bC$ to $\bB$. For every $x \in V$, define $f(x)=e([x]_\sim)$. Note that for every summand $R(x_1, \dots, x_k)$ in $\phi$ and $t \in f^*(y(x_1, \dots, x_k, \dots, x_k))$, we have
\begin{align*}
R^*(f^*(y(x_1, \dots, x_k, \dots, x_k)))  = R(t_1, \dots, t_k) = R(e([x_1]_\sim), \dots, e([x_k]_\sim)) = R(f(x_1), \dots, f(x_k)),
\end{align*}
where the middle equality follows from $t_i \mapsto e([x_i]_\sim)$ being a partial isomorphism of $\bB$ by Claim 1 and 2, which by the homogeneity of $\bB$ extends to an automorphism of $\bB$ and therefore also an automorphism of $\Gamma$.
Since $f^*$ is a solution to $(\phi^*,u)$, it follows from the construction of $\phi^*$ that $f$ is a solution to $(\phi,u)$.
\end{proof}

\end{document}
